% 可复制的空白模板。编译命令：./build.sh template
\documentclass{ai4s-resume}
% 改配色 / 页边距 / 字体大小时编辑本文件，不必动正文。
\definecolor{Accent}{HTML}{24455C} % 墨蓝；黑白版可改为 222222
% 字号与节奏。只改这些参数即可调整整份简历。
\renewcommand{\ResumeBodySize}{9.8}
\renewcommand{\ResumeLeading}{14.5}
\renewcommand{\ResumeSectionGap}{11pt}
\renewcommand{\ResumeEntryGap}{6pt}
\renewcommand{\ResumeProjectGap}{10pt}
% \geometry{left=15mm,right=15mm,top=14mm,bottom=16mm}
% 显示原 PDF 中的照片；关闭后页眉自动扩展为全宽。
\ResumePhotoEnabledtrue
% 两页版在“科研项目”前分页，避免长项目的关键指标被拆开。
% 内容增删较多时，可改为 false，让 LaTeX 自然分页。
\ResumeManualBreaktrue

\ResumePhotoEnabledfalse
\renewcommand{\ResumeName}{你的姓名}
\renewcommand{\ResumeTagline}{研究方向一 · 研究方向二 · 求职方向}
\renewcommand{\ResumePhone}{13800000000}
\renewcommand{\ResumeEmail}{your.name@example.com}
\renewcommand{\ResumeCity}{上海}
\renewcommand{\ResumeFooter}{研究方向 · 求职方向}
\begin{document}
\makeheader
\cvsection{教育经历}{EDUCATION}
\datedline{大学名称｜学院名称}{\dateRange{2024.09}{2027.06}}
\entrydetail{专业方向 · 实验室 · 硕士（在读）}
GPA / 排名；代表性奖项；与目标岗位相关的课程。
\cvsection{科研论文}{PUBLICATIONS}
\paper{Your Paper Title: A Clear Description of the Research Contribution.}{会议年份 · 真实状态 · 作者顺位}
\cvsection{实习经历}{EXPERIENCE}
\datedline{公司名称｜团队名称}{\dateRange{2025.06}{至今}}
\entrydetail{职位名称 · 研究方向}
\begin{itemize}
  \item 说明负责的具体问题、个人贡献和最终交付；按真实口径填写数据规模与实验结果。
  \item 给出评测数据集、对照方法、指标和实验设置，区分个人负责与参与的工作。
\end{itemize}
\cvsection{科研与项目经历}{RESEARCH \& PROJECTS}
\projecttitle{项目名称｜解决的问题}{所属团队 · 方法 / 技术路线}
\begin{itemize}
  \item \textbf{问题与方法：}描述任务、主要方法以及你负责的部分。
  \item \textbf{工程实现：}描述可复用系统、训练流程或评测框架。
  \item \textbf{实验结果：}在指定评测集上，指标从 \metric{基线数值} 到 \metric{最终数值}；注明统计口径。
\end{itemize}
\cvsection{专业技能}{TECHNICAL SKILLS}
\skill{训练与推理}{填写实际使用过的框架与分布式训练经验。}
\skill{研究方法}{填写后训练、多模态、Agent 或其他相关方法。}
\skill{工程能力}{填写编程语言、实验平台与工程工具。}
\end{document}
